
quelconques @@M0@@ et @@M1@@ sont conjuguées, autrement dit il existe une matrice inversible @@M2@@ @@M3@@ telle que l'on ait
@@M4@@
En effet, la matrice @@M5@@ avec @@M6@@, @@M7@@, et du reste @@M8@@, répond à la question (rappelons que @@M9@@ et @@M10@@ sont par hypothèse inversibles et @@M11@@). Donc pour un homomorphisme quelconque @@M12@@ de @@M13@@ dans un faisceau de groupes abéliens @@M14@@, on a
@@M15@@
Si @@M16@@ est inversible, on tire de ceci et de (6.3.1) que @@M17@@ pour tout @@M18@@, de sorte que @@M19@@ est un produit de commutateurs. Si le corps résiduel @@M20@@ de @@M21@@ n'est pas le corps premier @@M22@@, alors on peut trouver pour tout @@M23@@ inversible un voisinage zariskien @@M24@@ de @@M25@@ et un élément inversible @@M26@@ tel que @@M27@@ soit inversible. Il en découle que @@M28@@ est un produit de commutateurs sur @@M29@@.

Si @@M30@@, alors @@M31@@ est un voisinage étale de @@M32@@, dont le corps résiduel est différent de @@M33@@. On a gagné.

\subsection*{6.5. Remarques}
\textbf{6.5.1.} Lorsque les corps résiduels de @@M34@@ sont tous non isomorphes à @@M35@@, la démonstration précédente montre que 6.2.2. reste valable pour le topos annelé des faisceaux \emph{zariskiens} sur @@M36@@. En revanche, déjà pour @@M37@@, @@M38@@ est différent du groupe des commutateurs de @@M39@@. L'ordre de ce dernier est @@M40@@ tandis que l'ordre de @@M41@@ est @@M42@@ (cf. [1]).

\textbf{6.5.2.} Sur un topos localement annelé quelconque @@M43@@ il n'est pas vrai en général que toute représentation @@M44@@ soit définie par une
